\chapter{Implementacja i testy na telefonie}

Od rozpoczęcia prac do ostatniej implementacji upłynęły 653 dni.
Pierwsze dwie instalacje wykonano na urządzeniach:
\begin{itemize}
	\item HTC - Desire Z,
	\item Sony - Xperia E.
\end{itemize}

Były to wersje algorytmu wykorzystującego jeszcze dopasowywanie 
wzorca w~ostatnim etapie celem weryfikacji odnalezionych cyfr.
Ze względu na słabe wyniki, które zmusiły do przeprojektowania
rozwiązania, nie wykonano wtedy żadnych szczegółowych pomiarów.

Na urządzeniu HTC osiągnięto wydajność rzędu 4-6 klatek na sekundę
w~trybie wyszukiwania frontu. Odnalezienie frontu i~przejście 
w~tryb odczytu numeru powodowało spadek częstotliwości odświeżania
do 0.77 klatki na sekundę (rysunek \ref{fig:imp_htc_v1}).

\begin{figure}[h!]
    \centering
    \includegraphics[width=0.9\textwidth]{img/imp_htc_v1}
    \caption{Wyniki pierwszej wersji programu na urządzeniu HTC}
    \label{fig:imp_htc_v1}
\end{figure}

Na urządzeniu SONY osiągnięto wydajność rzędu 6-8 klatek na sekundę
w~trybie wyszukiwania frontu. Odczyt numeru był na poziomie 0.97
klatki na sekundę (rysunek \ref{fig:imp_sony_v1}).

\begin{figure}[h!]
    \centering
    \includegraphics[width=0.9\textwidth]{img/imp_sony_v1}
    \caption{Wyniki pierwszej wersji programu na urządzeniu SONY}
    \label{fig:imp_sony_v1}
\end{figure}

Trzeba pamiętać, że rozdzielczość na obu urządzeniach była ograniczona
rozdzielczością wyświetlacza \footnote{Biblioteka OpenCV do ciągłego pobierania
	danych z~aparatu urządzenia wykorzystuję funkcję podglądu 
	dostępną w~API systemu Android. Minusem takiego podejścia 
	jest zaniżona rozdzielczość pobieranych obrazów, która w~trybie podglądu jest ograniczona
	rozdzielczością wyświetlacza.}, która w~obu przypadkach była zbliżona:
\begin{itemize}
	\item HTC - 800x480,
	\item Sony - 768x432.
\end{itemize}

Będąc w~posiadaniu końcowej wersji aplikacji opartej na dwunastu
detektorach kaskadowych przystąpiono do szczegółowych pomiarów. 
Tym razem urządzeniem testowym był aparat LG G3. Większa rozdzielczość
wyświetlacza umożliwiała wykonywanie testów nawet w~rozdzielczości 1920x1440.
Biorąc pod uwagę ponad siedmiokrotne zwiększenie liczby pikseli 
w~obrazie do przeszukania, pierwszym etapem było zbadanie 
wydajności urządzenia na biegu jałowym. Wyniki pomiaru 
pobierania klatek z~aparatu przedstawiono w~następnym podrozdziale. 

\section{Testy wydajności podstawowych funkcji}

Podczas obcowania z~urządzeniem zaobserwowany został niepokojący fakt - telefon
będąc dłużej pod pełnym obciążeniem znacząco się rozgrzewał. Obudowa 
stawała się odczuwalnie gorąca. Podczas pierwszego kilkuminutowego
testu zbadano wpływ temperatury urządzenia na jego wydajność. Wyniki zaprezentowane
zostały na wykresie \ref{chart:time_2_temp_2_fps}. 

\begin{figure}[h!]
	\begin{center}
		\begin{tikzpicture}
		\begin{axis}[
		ylabel={$Liczba\ klatek\ na\ sekunde$},
		xlabel={$Liczba\ pobranych\ klatek$},
		minor y tick num=1,
		legend style={at={(0.05,0.30)},anchor=north west},
		width=\textwidth*0.9, height=6cm,
		/pgf/number format/.cd, use comma, 1000 sep={}
		]
		\addplot [only marks, color=blue] table {data/phone/raw1920x1440x01.dat};
		\addlegendentry{pierwsza proba}
		\addplot [only marks, color=red] table {data/phone/raw1920x1440x02.dat};
		\addlegendentry{druga proba}
		\end{axis}
		\end{tikzpicture}
	\end{center}
	\caption{Częstość pobierania klatek obrazu w~zależności od czasu (podczas 
		ciągłej pracy urządzenie stawało się odczuwalnie gorące). 
		Druga próba wykonana była po wystygnięciu urządzenia zaraz po pierwszej.
		Najwidoczniej urządzenie nie wystygło w~pełni i~szybciej się rozgrzało
		co spowodowało przejście urządzenie w~stan ograniczonej wydajności.}
	\label{chart:time_2_temp_2_fps}
\end{figure}

Zakładając, że w~implementacji bazowej dostarczonej wraz z~biblioteką 
OpenCV nie ma wycieków pamięci, można wysnuć tezę, że dwukrotny spadek
wydajności urządzenia spowodowany jest owym wzrostem temperatury.
Celem tego eksperymentu nie było wyznaczanie czasu
po którym następuje spadek wydajności, czy innych precyzyjnych parametrów.
Najważniejszym był fakt, że problem w~ogóle się pojawia i~trzeba brać go pod uwagę.
Może on nie dotyczyć wszystkich urządzeń lub objawiać się w~większym bądź
mniejszym stopniu. Ważnym jest aby podczas projektowania rozwiązania
mieć świadomość, że taki przypadek może mieć miejsce.
 Można zastosować wiele technik ograniczających zużycie 
procesora jak chociażby wprowadzenie oszczędnego trybu pracy podczas
wyszukiwania samego frontu. W~takim trybie aplikacja pobierałaby na przykład
co dziesiątą klatkę obrazu lub uruchamiała detektor
co sekundę. Powinno to znacząco ograniczyć wykorzystanie zasobów
obliczeniowych urządzenia.

Kolejnym testem było sprawdzenie wpływu wielkości obrazu na 
wydajność. Po ponad siedmiokrotnym zmniejszeniu rozdzielczości - z 1920x1440 na 
800x480 (jak na urządzeniach poprzedniej generacji: SONY i HTC) nie uzyskano 
zbyt wielkiego skoku wydajnościowego. 

\begin{figure}[h!]
	\begin{center}
		\begin{tikzpicture}
		\begin{axis}[
		ylabel={$Liczba\ klatek\ na\ sekunde$},
		xlabel={$Liczba\ pobranych\ klatek$},
		minor y tick num=1,
		legend style={at={(0.05,0.30)},anchor=north west},
		width=\textwidth*0.9, height=6cm,
		/pgf/number format/.cd, use comma, 1000 sep={}
		]
		\addplot [only marks, color=blue] table {data/phone/raw1920x1440x01.dat};
		\addlegendentry{rozdzielczosc 1920x1440}
		\addplot [only marks, color=green] table {data/phone/raw800x480x01.dat};
		\addlegendentry{rozdzielczosc 800x480}
		\end{axis}
		\end{tikzpicture}
	\end{center}
	\caption{Częstość pobierania klatek obrazu w~zależności od czasu oraz
		rozmiarów zdjęć pobieranych z~aparatu urządzenia}
	\label{chart:time_2_temp_2_fps_with_resolution_change}
\end{figure}

Jedyne co można zaobserwować na wykresie \ref{chart:time_2_temp_2_fps_with_resolution_change} to mniejsza podatność
na wzrost temperatury urządzenia, czyli mniejszy spadek wydajności
po dłuższej pracy pod obciążeniem.

\subsection{Pomiar czasu detekcji frontu na zdjęciu przeskalowanym do zadanej
	wysokości}

Przeskalowanie klatki obrazu o~wymiarach 1920x1440 pikseli do rozmiarów
300x225 wraz z~uruchomieniem detektora frontów w~tak przygotowanym obrazie
zajęło około 50 milisekund. Zdarzały się przypadki na poziomie poniżej
25 ms jednak wraz ze wzrostem temperatury urządzenia czas wykrywania 
frontu wzrastał, co zostało przedstawione na wykresie \ref{chart:detector_ms_time}.
Można tu odnotować trzy interesujące przedziały czasowe, dla których
wyliczono następujące średnie:
\begin{itemize}
	\item 1 - 1133 gdy urządzenie było jeszcze zupełnie zimne, średni czas
	wykonania: 44 ms,
	\item 1133 - 2611 zauważalny wzrost temperatury (pojedyncze przypadki na poziomie lub poniżej 50 ms), średni czas wykonania: 73 ms,
	\item 2612 - 3000 - przejście urządzenia w tryb oszczędny ze względu na temperaturę, średni czas wykonania: 131 ms,
	niemal wszystkie odczyty powyżej 100 ms.
\end{itemize}

\begin{figure}[h!]
	\begin{center}
		\begin{tikzpicture}
		\begin{axis}[
		ylabel={$Czas\ pobrania\ obrazu$},
		xlabel={$Liczba\ pobranych\ klatek$},
		minor y tick num=1,
		legend style={at={(0.05,0.90)},anchor=north west},
		width=\textwidth*0.9, height=6cm,
		/pgf/number format/.cd, use comma, 1000 sep={}
		]
		\addplot[mark=none, green, line width=2pt] coordinates {(0,44) (1150,44)};
		\addlegendentry{średni czas 44 ms}
		\addplot[mark=none, yellow, line width=2pt] coordinates {(1160,73) (2620,73)};
		\addlegendentry{średni czas 73 ms}
		\addplot[mark=none, red, line width=2pt] coordinates {(2630,131) (3000,131)};
		\addlegendentry{średni czas 131 ms}
		\addplot [only marks, mark size=0.5, color=green] table {data/phone/detector01.dat};
		\addplot [only marks, mark size=0.5, color=yellow] table {data/phone/detector02.dat};
		\addplot [only marks, mark size=0.5, color=red] table {data/phone/detector03.dat};
		\end{axis}
		\end{tikzpicture}
	\end{center}
	\caption{Rosnący czas, potrzebny na pobranie obrazu z~aparatu urządzenia, w~zależności
		od czasu pracy pod obciążeniem z~uwzględnieniem etapów kiedy urządzenie jest:
		zimne (kolor zielony), rozgrzewa się (kolor żółty), ze względu na 
		temperaturę zostało w~prowadzone w~tryb
		ograniczonej wydajności (kolor czerwony)}
	\label{chart:detector_ms_time}
\end{figure}

Im dłużej urządzenie pracowało na pełnych obrotach, tym bardziej wzrastała jego 
temperatura i~spadała wydajność. Co ważne, do momentu pobrania
nieco ponad tysiąca klatek przeskalowanie obrazu i~przeszukanie go
pod kątem wystąpienia frontu autobusu trwało nie więcej niż 75 milisekund,
a~znaczna większość przypadków trwała krócej niż 50 ms. Jest to bardzo 
ważna wartość progowa. Tyle właśnie czasu trwa pobranie pojedynczego
obrazu z~aparatu urządzenia bez żadnych dodatkowych modyfikacji.
Rozpatrując potencjalne zrównoleglenie procesu, czasy
pobrania klatki i~wyszukania frontu będące na tym samym poziomie 
umożliwiają optymalne wykorzystanie co najmniej dwóch rdzeni procesora urządzenia. 
Urządzenie testowe - LG G3 - wyposażone zostało w~czterordzeniowy procesor. Istnieje 
więc realna szansa wykonania całego procesu nawet do~20 razy na sekundę. 
Jednak wszystkie dotychczasowe wersje testowe jak i~program końcowy zaprojektowane i~zaimplementowane
zostały w~sposób jednowątkowy. Osiągnięcie wydajności na poziomie 20 odczytów na
sekundę jest więc jednym z~proponowanych usprawnień.

Wykrywanie numeru okazało się zadaniem dużo mniej złożonym obliczeniowo.
Średni czas wykrycia numeru we froncie to 16 ms. Założenie, że nigdy nie
przekroczy on 40 ms jest bardziej niż pewne.
To, z~kolei ,
umożliwia przeszukanie 25 frontów pod kątem wystąpienia numeru na sekundę.
W~razie potrzeby proces ten mógłby być obsługiwany przez kolejny wątek.
Wąskim gardłem byłyby w~tym przypadku operacje pobrania obrazu 
z~aparatu urządzenia oraz przeskalowania i~przeszukana go pod kątem wystąpienia 
frontu autobusu. Obie operacje, nawet wykonane równolegle uniemożliwiają 
wykonanie całego algorytmu częściej niż 20 razy na sekundę.

\begin{figure}[h!]
	\begin{center}
		\begin{tikzpicture}
		\begin{axis}[
		ylabel={$Czas\ wykonania\ [ms]$},
		xlabel absolute, xlabel style={yshift=-0.5cm},
		xlabel={$Liczba\ pobranych\ klatek$},
		minor y tick num=1,
		x tick label style={rotate=45},
		legend style={at={(0.05,0.90)},anchor=north west},
		width=\textwidth*0.9, height=6cm,
		/pgf/number format/.cd, use comma, 1000 sep={}
		]
		\addplot [only marks, color=blue] table {data/phone/digits01.dat};
		%\addlegendentry{czas odczytu numeru linii}
		\end{axis}
		\end{tikzpicture}
	\end{center}
	\caption{Czas potrzebny na odczytanie numeru linii z~fragmentu
		obrazu go reprezentującego w~zależności liczby pobranych 
		i~przetworzonych
		klatek}
	\label{chart:digits_ms_time}
\end{figure}

Odczyt wartości widocznej na obrazie reprezentującym numer linii
trwał średnio 18 ms. Była to wartość uzyskana podczas pracy
na zimnym urządzeniu. Po dokonaniu ponad 1900 odczytów
wartość średniego czasu wykonania wzrosła do 40 ms.
Moment przejścia urządzenia w~tryb oszczędny (chroniący go przed
osiągnięciem nadmiernej temperatury) jest dobrze widoczny na wykresie
\ref{chart:digits_ms_time}.

W~ramach podsumowania przygotowano zbiorcze wykresy 
przedstawiające sumaryczny czas wykonania wszystkich wspomnianych
etapów:
\begin{enumerate}
    \item Pobrania obrazu z~aparatu urządzenia.
    \item Przeskalowania klatki do 300 pikseli wysokości i~przeszukania 
    jej pod kątem wystąpienia frontu autobusu.
    \item Przeskalowania wycinka przedstawiającego front do 250 pikseli wysokości
    i~przeszukania go pod kątem wystąpienia numeru (w~lewym górnym rogu obrazu).
    \item Przeskalowanie wycinka przedstawiającego numer do 50 pikseli wysokości
    i~przeszukania go z~wykorzystaniem dziesięciu detektorów cyfr oraz~przygotowania
    ciągu znaków reprezentującego numer.
\end{enumerate}

Trzy wykresy reprezentują trzy wspomniane stany urządzenia: zimne,
ciepłe (przed przejściem w~tryb oszczędny), gorące (w~trybie oszczędnym).

\begin{figure}[h!]
	\begin{center}
		\begin{tikzpicture}
		\begin{axis}[
		ylabel={$Czas\ wykonania\ [ms]$},
		xlabel absolute,
		xlabel={$Liczba\ pobranych\ klatek$},
		stack plots=y,
		area style,      
		ymin=0,         
		minor y tick num=1,   
		enlarge x limits=false,
		enlarge y limits=upper,
		width=\textwidth*0.9, height=6cm,
		legend style={at={(0.575,0.50)},anchor=north west},
		/pgf/number format/.cd, use comma, 1000 sep={}
		]
		\addplot+[y] table {data/phone/allxother.dat} \closedcycle;
		\addlegendentry{pobranie obrazu z aparatu}
		\addplot+[y] table {data/phone/allxfront.dat} \closedcycle ;
		\addlegendentry{wyszukanie frontu}
		\addplot+[y] table {data/phone/allxnumber.dat} \closedcycle;
		\addlegendentry{wyszukanie numeru}
		\addplot+[y] table {data/phone/allxdigits.dat} \closedcycle;
		\addlegendentry{odczytanie numeru}
		
		\end{axis}
		\end{tikzpicture}
	\end{center}
	\caption{Ogólny czas wykonania całego algorytmu na zimnym urządzeniu}
	\label{chart:whole_alg_ms_time_cold}
\end{figure}

Gdy urządzenie było zupełnie zimne - pierwsze 600 pobranych klatek - 
ogólny czas wykonania nie przekraczał 150 ms (dane widoczne na wykresie
\ref{chart:whole_alg_ms_time_cold}). Wyodrębnić można podział na 
maksymalnie trzy wątki (wyszukanie i~odczytanie numeru jako jeden wątek), 
co w~rezultacie umożliwiłoby wykonanie pełnego cyklu w~czasie równym
pobraniu obrazu z~aparatu: 50 ms (20 klatek na sekundę).

\begin{figure}[h!]
	\begin{center}
		\begin{tikzpicture}
		\begin{axis}[
		ylabel={$Czas\ wykonania\ [ms]$},
		xlabel absolute,
		xlabel={$Liczba\ pobranych\ klatek$},
		stack plots=y,
		area style,      
		ymin=0,         
		minor y tick num=1,   
		enlarge x limits=false,
		enlarge y limits=upper,
		width=\textwidth*0.9, height=6cm,
		legend style={at={(0.575,0.50)},anchor=north west},
		/pgf/number format/.cd, use comma, 1000 sep={}
		]
		\addplot+[y] table {data/phone/allxother02.dat} \closedcycle;
		\addlegendentry{pobranie obrazu z aparatu}
		\addplot+[y] table {data/phone/allxfront02.dat} \closedcycle ;
		\addlegendentry{wyszukanie frontu}
		\addplot+[y] table {data/phone/allxnumber02.dat} \closedcycle;
		\addlegendentry{wyszukanie numeru}
		\addplot+[y] table {data/phone/allxdigits02.dat} \closedcycle;
		\addlegendentry{odczytanie numeru}
		\end{axis}
		\end{tikzpicture}
	\end{center}
	\caption{Ogólny czas wykonania całego algorytmu na ciepłym urządzeniu}
	\label{chart:whole_alg_ms_time_warm}
\end{figure}

Wzrost temperatury skutkował pojedynczymi wydłużonymi czasami wykonania,
które trwały ponad 200 ms (dane zaprezentowane na wykresie
\ref{chart:whole_alg_ms_time_warm}). Urządzenie nie przeszło wtedy
jeszcze w~tryb zmniejszonej częstotliwości procesora.

\begin{figure}[h!]
	\begin{center}
		\begin{tikzpicture}
		\begin{axis}[
		ylabel={$Czas\ wykonania\ [ms]$},
		xlabel absolute,
		xlabel={$Liczba\ pobranych\ klatek$},
		stack plots=y,
		area style,      
		ymin=0,         
		minor y tick num=1,   
		enlarge x limits=false,
		enlarge y limits=upper,
		width=\textwidth*0.9, height=6cm,
		legend style={at={(0.575,0.50)},anchor=north west},
		/pgf/number format/.cd, use comma, 1000 sep={}
		]
		\addplot+[y] table {data/phone/allxother03.dat} \closedcycle;
		\addlegendentry{pobranie obrazu z aparatu}
		\addplot+[y] table {data/phone/allxfront03.dat} \closedcycle ;
		\addlegendentry{wyszukanie frontu}
		\addplot+[y] table {data/phone/allxnumber03.dat} \closedcycle;
		\addlegendentry{wyszukanie numeru}
		\addplot+[y] table {data/phone/allxdigits03.dat} \closedcycle;
		\addlegendentry{odczytanie numeru}
		\end{axis}
		\end{tikzpicture}
	\end{center}
	\caption{Ogólny czas wykonania całego algorytmu na gorącym urządzniu
		pracującym w~trybie oszczędnym}
	\label{chart:whole_alg_ms_time_hot}
\end{figure}

Po dostatecznie długiej pracy i~pobraniu ponad 1600 klatek obrazu urządzenie
przeszło w~tryb oszczędny. Wydajność spadła prawie dokładnie dwukrotnie,
a~czasy wykonania całego algorytmu wydłużyły się do niemal 300 milisekund
(wykres \ref{chart:whole_alg_ms_time_hot}). Jest to problem, który
na pewno trzeba by wziąć pod uwagę podczas projektowania
i~implementacji komercyjnego rozwiązania.
Na potrzeby niniejszej pracy wystarczy mieć świadomość, że zjawisko
występuje. Przy założeniu, że aplikacja obsługiwana byłaby przez 
trzy wątki oraz po wprowadzeniu urządzenia w~tryb ograniczonej wydajności
czas wykonania pojedynczego odczytu wzrósłby do 100 milisekund. Czyli nadal
możliwe byłoby pobranie i~przetworzenie 10 klatek obrazu w~ciągu jednej sekundy.
Zakładając, że od wykrycia autobusu do zatrzymania jest około 2-4 sekundy 
nadal dysponujemy wystarczającą liczbą obrazów do wykonania odczytu numeru.
Problemem byłaby zatem temperatura urządzenia sama w~sobie oraz
drastycznie skrócony czas pracy telefonu na baterii.
	
\subsection{Testy skuteczności na urządzeniu}

Do przetestowania skuteczności wybrano zbiór 3100 wycinków
obrazów reprezentujących fronty autobusów z~wyświetlaczem
diodowym. Pojedynczy obraz wyświetlany był wielokrotnie
przy jednoczesnym zwiększaniu jego rozmiarów.
Począwszy od kwadratu o~rozmiarach 500x500 pikseli, 
każdorazowo zwiększano wysokość i~szerokość wyświetlanego obrazu
o~10 pikseli. Po pięćdziesięciu próbach odczytu (oraz osiągnięciu
wymiarów równych 1000x1000 pikseli) cała procedura była powtarzana
dla kolejnego obrazu.
Celem tego ćwiczenia była jak najwierniejsza
 symulacja sytuacji z~przystanku autobusowego.
Wyświetlane w~ten sposób obrazy były odczytywane przez
telefon ustawiony na statywie przed monitorem.
Dane z~telefonu przekazywane były do komputera poprzez 
kabel USB, a~następnie porównywane z~wartościami określającymi
numer linii, które zawarte były w~nazwach plików.

Taki wstępny test tylko
na pozór symulował sytuację z~przystanku. Nie do odtworzenia
była bowiem zmiana oświetlenia i~geometrii, a~co za tym idzie
stopnia trudności spowodowanego przez zakłócenia. Teoretycznie,
wykorzystując proponowaną metodę, osiągnięte wyniki
powinny być znacznie gorsze niż w~rzeczywistości. 
Jeżeli bowiem oryginalne zdjęcie zawierało odblaski zupełnie 
uniemożliwiające odczytanie numeru, to stopniowe zwiększanie obrazu 
nie usunęło tego problemu.
Numer nie został odczytany. W~naturalnych warunkach odblaski
na szybie powinny się ciągle zmieniać, przez co prawdopodobieństwo 
poprawnego odczytu powinno być dużo większe. Niemożliwa byłaby bowiem sytuacja,
 aby wraz z~przemieszczaniem się autobusu
promienie słoneczne odbijały się cały czas w~sposób uniemożliwiający odczyt.

Na wykresie \ref{chart:final_device_test} przedstawiono wyniki testu.
Niebieskie kropki reprezentują liczbę obrazów w~poszczególnych
grupach skuteczności. Jeżeli podczas prób z~danym obrazem, wartość
poprawnych odczytów była równa 50 (czyli maksymalnej liczbie wyświetlanych obrazów), obraz taki umieszczany był w~katalogu o~nazwie 100. Suma wszystkich
obrazów, per katalog, została zaprezentowana na wspomnianym wykresie
numer \ref{chart:final_device_test}. 

\begin{figure}[h!]
	\begin{center}
		\begin{tikzpicture}
		\begin{axis}[
		ylabel={$Liczba\ obrazow\ w\ danej\ grupie$},
		xlabel={$Grupy\ procentowej\ skutecznosci\ detekcji\ calego\ systemu$},
		minor y tick num=1,
		legend style={at={(0.05,0.30)},anchor=north west},
		width=\textwidth*0.9, height=6cm
		]
		\addplot [only marks, color=blue] table {data/final/result.dat};
		%\addlegendentry{pierwsza proba}
		\end{axis}
		\end{tikzpicture}
	\end{center}
	\caption{Reprezentacja podzbiorów obrazów, dla których
		testowany system osiągnął daną skuteczność wyrażoną w~procentach 
		(oś x) oraz liczba elementów przypisanych do danej grupy (oś y)}
	\label{chart:final_device_test}
\end{figure}

Najliczniejszą grupą - niemal 300 - 
są obrazy dla których nie odnotowano ani jednego poprawnego odczytu.
Kolejna liczna grupa - nieco ponad 200 - to próbki dla których
co najmniej jeden odczyt był już prawidłowy. Co ciekawe i~ważne, odnotowano 98
przypadków dla których wszystkie 50 odczytów było poprawnych.

Patrząc na powyższy wykres można stwierdzić, że rozwiązanie jest 
dalekie od doskonałości. Sposób mierzenia skuteczności 
był również bardzo niedoskonały i~niepoprawnie modelował rzeczywistość.
W~nie-laboratoryjnych warunkach wyniki powinny być raczej dużo lepsze.

Ważnym produktem ubocznym omawianego doświadczenia był zestaw podzbiorów obrazów
podzielonych i~posortowanych zgodnie z~rosnącą trudnością ich poprawnej obsługi.
Posiadając tak spreparowane zbiory, można było w~łatwy sposób zweryfikować
jakie były słabe elementy systemu oraz jakie przypadki użycia
sprawiały aplikacji największe problemy.

\subsection{Ręczna weryfikacja wyników na urządzeniu}

Po przygotowaniu podzbiorów, posortowanych zgodnie z~rosnącą
trudnością w~ich odczytaniu, przystąpiono do ręcznej weryfikacji wyników.
Zupełnie pominięte zostały fronty autobusów z~odblaskami widocznymi
w~szybie osłaniającej wyświetlacz. Jak już to zostało 
wspomniane problem w~warunkach naturalnych nie powinien występować,
a~przynajmniej nie w~tej skali.

Kolejnym problemem był dobór samych frontów. Pominąwszy
pojedyncze egzemplarze w~których numer był umieszczony z~prawej
strony - niespełnienie założenia wstępnego - skupiono się na wychwyceniu nietypowych
przypadków. Większość eksperymentów odbywała się na zbiorze obrazów,
dla których nie odnotowano ani jednego poprawnego odczytu.

\begin{figure}[h!]
	\centering
	\includegraphics[width=1\textwidth]{img/final/front}
	\caption{Pięć przykładowych frontów, które nie zostały wykryte
		podczas automatycznej sesji testowej - na czarnym 
		pasku umieszczono wyniki trzech ręcznie wykonanych prób
		weryfikacyjnych}
	\label{fig:final_front}
\end{figure}

Na rysunku \ref{fig:final_front} 
przedstawionych został pięć przykładowych frontów, których
kształt nie pozwalał na ich lokalizację. Brak lokalizacji frontu
uniemożliwiał dalsze działanie algorytmu. Kontrprzykład, czyli autobus
linii 174, który udało się odczytać podczas manualnej sesji
testowej, unaocznia pewne niedociągnięcia
przygotowanej procedury testowej. Zakładała ona jedynie 
przeskalowanie, czyli ustawiczne powiększanie 
obrazu reprezentującego front. Nie uwzględnione były inne 
przekształcenia geometryczne, perspektywiczne i~różnice w~oświetleniu.
W~tym konkretnym przypadku po niedługim manipulowaniu aparatem udało się
odnaleźć właściwy kąt dla którego odczyt przebiegał bez zakłóceń.

Kolejny problem to utrwalone na zdjęciach zniekształcenia geometryczne 
prezentowanych frontów.
Także w~tym przypadku,
udało się poprawić wyniki podczas ręcznego operowania telefonem.
Przykłady zamieszczono na obrazie nr \ref{fig:final_deform}.

\begin{figure}[h!]
	\centering
	\includegraphics[width=1\textwidth]{img/final/deform}
	\caption{Fronty autobusów, które ze względu na duże zniekształcenia 
		geometryczne nie zostały wykryte przez detektor frontów -
		na czarnym pasku zamieszczono wyniki trzech ręcznych prób
		wykonanych w~celu weryfikacji testu automatycznego}
	\label{fig:final_deform}
\end{figure}

Po odpowiednim nachyleniu aparatu fronty z~pierwszych dwóch obrazów
zostały na tyle ,,wyprostowane'', że udało się uzyskać poprawny odczyt.

O~ile poprzednia klasa problemów, wymagałaby przeprojektowania 
rozwiązania lub co najmniej przygotowania osobnego zestawu
detektorów frontów, to zjawisko zniekształceń w~praktyce
w~ogóle nie powinno się pojawić.
Stojący na przystanku obserwator niemal zawsze będzie widział
front autobusu idealnie zwrócony w~jego stronę. Zniekształcenia
geometryczne mogą pojawić się dopiero przy bardzo małych odległościach
od obserwowanego obiektu.

\begin{figure}[h!]
	\centering
	\includegraphics[width=1\textwidth]{img/final/jednocy}
	\caption{Przykładowy podzbiór autobusów z~numerem jednocyfrowym.
		Dodatkowym utrudnieniem w~dwóch pierwszych przykładach była
		niewielka odległość między numerami, a~tekstem opisującym trasę.}
	\label{fig:final_jednocy}
\end{figure}

Pierwsze znaczące błędy
i~praktyczny brak odczytów pojawił się dla liczb jednocyfrowych.
Rozwiązanie od początku było projektowane z~myślą o~numerach trzycyfrowych,
które jako jedyne występowały w~środowisku testowym (miasto stołeczne).
Numery jednocyfrowe stanowią najtrudniejszą do rozwiązania klasę
problemów. Po pierwsze są one niechętnie wykrywane przez detektor
numerów, który został przygotowany głównie na zbiorach trzycyfrowych.
Po drugie, przy niewielkich odległościach numeru od opisu trasy,
pojawia się niebezpieczeństwo potraktowania pierwszych liter opisu jako cyfr.
Można wtedy błędnie uzyskać numer trzycyfrowy który nijak się ma do
faktycznego numeru linii autobusowej. Przykład
takiego zachowania przedstawiono na rysunku \ref{fig:final_jednocy}.

\begin{figure}[h!]
	\centering
	\includegraphics[width=1\textwidth]{img/final/waskie}
	\caption{Przykłady autobusów z~wyświetlaczami wykorzystującymi 
		wąskie czcionki. Na czarnych paskach przedstawiono wyniki
		trzech manualnych odczytów testowych.}
	\label{fig:final_waskie}
\end{figure}

Kolejne dwa przypadki to tak naprawdę dwa aspekty tego samego zagadnienia.
Wynikały one z~faktu, że przyjęta w~programie odległość między poszczególnymi 
cyframi nie była elastyczna i~ustawiono ją na konkretną wartość.
Testy dla pierwszych pięciu
sekwencji przygotowanych na podstawie filmów nakręconych na przystankach
wykazały, że optymalna odległość między cyframi, dla obszaru obrazu
reprezentującego numer przeskalowanego do 250 pikseli wysokości to 10 pikseli.
Był to pierwszy element algorytmu, którego implementację należało jeszcze
raz przemyśleć i~udoskonalić.
Kilka przykładów, wraz z~odczytanymi wartościami 
zaprezentowano na rysunku \ref{fig:final_waskie}, 
gdzie czcionki były nieco węższe niż te które posłużyły jako pierwowzór.

\begin{figure}[h!]
	\centering
	\includegraphics[width=1\textwidth]{img/final/szerokie}
	\caption{Przykłady autobusów z~wyświetlaczami wykorzystującymi 
		szerokie czcionki. Na czarnych paskach przedstawiono wyniki
		trzech manualnych odczytów testowych.}
	\label{fig:final_szerokie}
\end{figure}

Z~drugiej strony gdy czcionki były szersze - jak na rysunku \ref{fig:final_szerokie} -
elementem uniemożliwiającym odczyt nie była z~góry zadana odległość
między cyframi, a~sam detektor numerów.
Obszar oznaczany jako numer był zbyt mały, w~wyniku czego do funkcji
odczytującej przekazywane były jedynie dwie pierwsze cyfry,
gdzie ostatnia była zupełnie pominięta.
Prostym rozwiązaniem tego problemu mogłoby być
dodanie kilku pikseli właśnie z~owej prawej strony, tak by obszar wynikowy
zawierał wszystkie trzy cyfry.

Krój czcionki oraz zmienna liczba cyfr w~numerze
nie zostały wzięte pod uwagę podczas projektowania i~implementacji. 
Szczególnie trudny problem stanowią numery jednocyfrowe,
z~niedalekim sąsiedztwem tekstu opisującego trasę. Wprowadzenie
usprawnienia na tym polu byłoby jednym z~ostatnich działań, które 
można poczynić w~celu poprawy jakości i~skuteczności proponowanego rozwiązania.
To i~kilka innych drobnych usprawnień, jak na przykład osobne zestawy
detektorów dla różnych typów frontów, czy weryfikacja możliwych
numerów na podstawie współrzędnych geograficznych, umożliwiłyby
uzyskanie skuteczności na zadowalającym poziomie.
Wykracza to jednak poza zakres tej pracy.

\subsection{Testy w terenie}

Wyniki z~poprzedniego podrozdziału nie 
przekładały się na to, co zaobserwowano podczas testów terenowych. Po wykonaniu 
kilkudziesięciu sesji testowych odnotowano tylko jeden niepoprawny odczyt.
Próby wykonywano tak, aby po zatrzymaniu autobus był nadal kilka metrów
przed obserwatorem. W~ten sposób, przy liczbie obrazów niezbędnych do wykonania
odczytu ustawionej na 3, do pełnego zatrzymania się autobusu 
możliwe było wykonanie dwóch, czasami nawet trzech odczytów. Dla takiej
konfiguracji wszystkie, z~wyjątkiem jednej, sesje były zakończone
w~stu procentach poprawnymi wynikami. Na przystanku Metro Ratusz Arsenał 
dla linii 227 w~sesji uzyskano trzy odczyty przy czym dwa
z~nich były poprawne i~jeden
niepoprawny (27).

Kolejną zidentyfikowaną słabą cechą proponowanej aplikacji była maksymalna odległość
autobusu umożliwiająca poprawny odczyt. Wynosiła ona od 
20 do 25 metrów. Ostatnią 
czynnością wykonaną w~ramach tej pracy były
próby zwiększenia tej wartości.

\begin{figure}[h!]
	\centering
	\includegraphics[width=0.6\textwidth]{img/final/framescaling/400height}
	\caption{Rozdzielczość obrazu przedstawiającego numer linii po
		przeskalowaniu klatki wejściowej do wysokości równej 400 pikseli 
		z~zachowaniem proporcji}
	\label{fig:final_resize_400}
\end{figure}


\subsection{Testy zasięgu}

Jako ostatni przeprowadzono test
ze zmiennym parametrem pierwszego elementu kaskady, czyli wysokością
do jakiej przeskalowywane były oryginalne klatki na potrzeby wykrywania
frontu autobusu.

W~pierwszej iteracji zwiększono wysokość z~300 do 400 pikseli. 
Przykładowe fragmenty obrazów reprezentujące numer, jakie były
pobierane dla najmniejszego frontu jaki można było zlokalizować
przedstawiono na rysunku \ref{fig:final_resize_400}.
Skuteczność ani wydajność nie uległy zmianie. Bardzo prawdopodobnym jest,
że zmiana ta nie będzie mieć negatywnego wpływu na ogólne działanie 
aplikacji. Zwiększy się natomiast tylko jej zasięg.

\begin{figure}[h!]
	\centering
	\includegraphics[width=0.6\textwidth]{img/final/framescaling/500height}
	\caption{Rozdzielczość obrazu przedstawiającego numer linii po
		przeskalowaniu klatki wejściowej do wysokości równej 500 pikseli 
		z~zachowaniem proporcji}
	\label{fig:final_resize_500}
\end{figure}

Drugi z~serii i~ostatni w~ogóle, był test wykonany dla wysokości
ustawionej na 500 pikseli. Przykładowe fragmenty obrazów
przedstawiające numer zamieszczono na rysunku \ref{fig:final_resize_500}.
Dla tej sesji testowej odczyty były już dużo mniej stabilne.
Średnio co drugi odczyt był niepoprawny, a~niewielkie ruchy aparatem
powodowały zupełny jego brak. Oczywiście, aby ustalić dokładną wartość trzeba
by przeprowadzić kolejne testy automatyczne, jednak ze względu na brak
rzetelnej procedury testowej nie zostały one przeprowadzone ponownie.
Sama zastosowana metoda wydaje się być na tyle niestabilna, że
jakiekolwiek określanie jej skuteczności na tym etapie było niepożądane.
Przed następną sesją testową rozsądnym byłoby usprawnienie
elementów, błędów oraz niedociągnięć, które zostały zlokalizowane, takich
jak niedokładność detektorów dla mniej niż trzycyfrowych numerów
lub problemy z~wyznaczaniem liczby cyfr dla wąskich czcionek. 

Aplikacja bez wprowadzania wyżej opisanych usprawnień udostępniała
minimalną funkcjonalność. Implementacja była daleka od optymalnej 
jednak przeprowadzone testy pozwoliły określić poziom skuteczności,
jakiego można się po niej w~obecnym stanie spodziewać. Zastosowana architektura
umożliwia poprawę jej działania poprzez dopracowywanie poszczególnych elementów,
dodawanie nowych lub wymianę mało skutecznych kroków na zupełnie nowe
rozwiązania. Duży potencjał w~kontekście odczytywania już zlokalizowanego
numeru drzemie w~sieciach neuronowych, jednak podobnie 
jak to miało miejsce w~przypadku detektorów tak i~do przygotowania
sieci neuronowych potrzebne byłyby ogromne zbiory danych i~takie same ilości czasu. 
Powstrzymano się więc od kolejnych usprawnień, udoskonaleń i~przeprowadzania
dodatkowych testów z~tym związanych. Do działającej, rzetelnej
i~budzącej zaufanie wersji aplikacji pozostała jeszcze bardzo długa droga.\footnote{Jedną z~uwag jakie otrzymałem podczas redagowania
	tekstu, był zbyt krytyczny stosunek do rzetelności przygotowanego
	oprogramowania. Jednak gdybym to ja był osobą niewidomą i~miał polegać
	tylko na odczytach oferowanych przez proponowaną aplikację,
	byłbym w~nie lada tarapatach. Skuteczność na poziomie 80\% to dużo,
	jednak jeden na pięć autobusów wiozących w~nieodpowiednie miejsce
	to już inna sprawa. Mówiąc o~wersji rzetelnej i~budzącej zaufanie,
	miałem na myśli gotowy produkt, który można by zaoferować użytkownikom
	końcowym. Produktem niniejszej pracy jest wczesna wersja prototypowa.}